% Pista pro-26j-q3 · Probabilitat
% Procedència: PAU Catalunya 2026 · Convocatòria de juny · Sèrie 1 · Exercici 3
\documentclass[11pt,a4paper]{article}
\usepackage{pista}
\pistainfo{Catalunya 2026}{Convocatòria de juny}{Sèrie 1}

\begin{document}

\section*{\textcolor{blauPista}{Exercici 3}}

\noindent L'Ajuntament de Canet de Mar ha aconseguit 24 entrades gratuïtes per a un concert d'un grup de rock català, i ha decidit sortejar-les entre els veïns interessats a assistir-hi. Tots els veïns del poble seguidors d'aquest grup s'apunten al sorteig, i només al $16\,\%$ els toca una entrada. De la resta de seguidors del grup, sis setenes parts intenten comprar una entrada per la web, on la probabilitat d'aconseguir-ne és del $25\,\%$.

\subsection*{3.1. \normalfont Quantes persones d'aquest municipi tenen entrada per al concert? \hfill\textnormal{[0,75 punts]}}

\pista{Calcula per separat les dues ``vies'' per aconseguir entrada: la probabilitat d'obtenir-la al sorteig ($16\,\%$) i la probabilitat d'aconseguir-la després per la web (compte: aquesta és una probabilitat composta, ja que primer cal no haver tingut sort al sorteig, després haver-ho intentat --sis setenes parts-- i finalment haver-ho aconseguit --el $25\,\%$--). Suma els dos percentatges per obtenir el percentatge total de seguidors amb entrada. Finalment, fes servir la dada que el $16\,\%$ correspon a 24 persones per plantejar una regla de tres i trobar quantes persones representa el percentatge total.}

\subsection*{3.2. \normalfont Si escollim a l'atzar un veí de Canet seguidor d'aquest grup de rock i no té entrada per al concert, quina és la probabilitat que hagués intentat aconseguir-la via web? \hfill\textnormal{[0,75 punts]}}

\pista{Et demanen una probabilitat condicionada ``en sentit invers'' respecte de les dades que et donen (que parlaven de no tenir èxit \emph{sabent que} s'ha intentat per la web): és un cas clar per aplicar el teorema de Bayes. Necessitaràs $P(\text{no té entrada})$ (que és el complementari de la probabilitat total que has calculat a l'apartat anterior) i $P(\text{no té entrada} \mid \text{ho intenta per la web})$ (que és el complementari del $25\,\%$ d'èxit al web).}

\vspace{0.6em}
\noindent El dia del concert, l'equip de so mesura el nivell de decibels generat pels crits entusiastes del públic durant els 5 minuts de durada de la cançó més famosa del grup; es pot aproximar per la funció
\[
S(t) = -t^{3} + 12t^{2} - 30t + 90, \qquad t \in [0, 5],
\]
on $t$ és el temps en minuts i $S(t)$ els decibels. Quan se superen els 100 decibels es considera que el públic està molt entregat i s'activen automàticament uns efectes lumínics especials.

\subsection*{3.3. \normalfont S'activaran en algun moment durant aquests cinc minuts? Si la resposta és afirmativa, calculeu en quin minut s'activen, aproximat a les dècimes. \hfill\textnormal{[1 punt]}}

\pista{Reformula la pregunta: vols saber si, en algun moment, $S(t) = 100$. Defineix la funció auxiliar $D(t) = S(t) - 100$ i comprova que és contínua (és polinòmica). Avalua $D(0)$ i $D(5)$: si tenen signes oposats, el teorema de Bolzano et garanteix que hi ha un punt on $D$ s'anul·la (és a dir, on $S(t) = 100$). Per aproximar aquest punt a les dècimes, ves avaluant $D$ en punts intermedis --per exemple, partint successivament per la meitat l'interval on saps que hi ha el canvi de signe-- fins que la diferència entre dos valors consecutius de $t$ sigui d'una dècima.}

\end{document}
